\chapter{VAT Module Implementation Specification}\label{ch:vat-impl}

This chapter provides the technical implementation specification for the VAT
processing module, addressing the review finding that ``VAT engine rule format
is underspecified.'' The specification enables developers to build conformant
implementations.

\section{Module architecture}

The VAT module comprises four components:

\begin{enumerate}
  \item \textbf{RulesEngine} --- evaluates JSON-Logic expressions against invoice
        data and produces checklist items;
  \item \textbf{RegistrationValidator} --- validates VAT registration numbers
        against country-specific formats and (optionally) authority APIs;
  \item \textbf{Calculator} --- computes expected VAT amounts given rates and
        subtotals;
  \item \textbf{ActivityLookup} --- maps business activity codes to VAT
        treatment categories.
\end{enumerate}

\begin{figure}[htbp]
\centering
\begin{tikzpicture}[
  node distance=8mm and 15mm,
  box/.style={draw,rounded corners,minimum width=28mm,minimum height=10mm,align=center,font=\small,fill=blue!5},
  db/.style={draw,cylinder,shape border rotate=90,aspect=0.25,minimum width=18mm,minimum height=12mm,fill=gray!10,font=\small},
  t/.style={-{Latex[length=2mm]}}
]
\node[box] (eng) {RulesEngine};
\node[box,below left=of eng,align=center] (reg) {Registration\\Validator};
\node[box,below=of eng] (calc) {Calculator};
\node[box,below right=of eng,align=center] (act) {Activity\\Lookup};
\node[db,above=of eng] (rules) {rules.json};
\node[db,above left=of eng,xshift=-5mm,align=center] (inv) {invoice\\record};
\node[db,above right=of eng,xshift=5mm,align=center] (out) {checklist\\record};

\draw[t] (inv.south) -- (eng.north west);
\draw[t] (rules.south) -- (eng.north);
\draw[t] (eng.north east) -- (out.south);
\draw[t,dashed] (eng.south west) -- (reg.north east);
\draw[t,dashed] (eng.south) -- (calc.north);
\draw[t,dashed] (eng.south east) -- (act.north west);
\end{tikzpicture}
\caption{VAT module component diagram. Dashed arrows indicate helper
         invocations from the rules engine.}
\label{fig:vat-arch}
\end{figure}

\section{API contracts}

Each helper function has a defined API contract. Implementations \textbf{must}
conform to these signatures to be invoked by the rules engine.

\subsection{Registration validation}

\begin{lstlisting}[language=Python,basicstyle=\ttfamily\small,caption={Registration validator API contract.}]
def check_vat_registration_validity(
    tax_id: str,
    country: str  # ISO 3166-1 alpha-2
) -> ValidationResult:
    """
    Validate a VAT registration number.
    
    Returns:
        ValidationResult with fields:
        - valid: bool
        - format_valid: bool
        - authority_confirmed: Optional[bool]
        - error_code: Optional[str]
    """
\end{lstlisting}

\begin{table}[htbp]
\centering
\caption{VAT registration number formats by country.}
\label{tab:vat-formats}
\footnotesize
\begin{tabularx}{\linewidth}{@{}l l X@{}}
\toprule
\textbf{Country} & \textbf{Regex} & \textbf{Notes} \\ \midrule
SA & \verb|^[0-9]{15}$| & 15-digit ZATCA TIN \\
EG & \verb|^[0-9]{9}$| & 9-digit ETA TRN \\
BH & \verb|^[0-9]{9,15}$| & NBR VAT number (variable length) \\
QA & \verb|^[A-Z0-9]{10}$| & QCSD registration (non-VAT) \\
\bottomrule
\end{tabularx}
\end{table}

\subsection{VAT calculation}

\begin{lstlisting}[language=Python,basicstyle=\ttfamily\small,caption={VAT calculator API contract.}]
def calculate_vat_amount(
    subtotal: Decimal,
    vat_rate_percent: Decimal
) -> CalculationResult:
    """
    Calculate expected VAT amount.
    
    Returns:
        CalculationResult with fields:
        - vat_amount: Decimal (rounded to 2 decimal places)
        - gross_amount: Decimal
        - calculation_method: str  # "standard" | "reverse"
    """
\end{lstlisting}

The calculator \textbf{must} use banker's rounding (ROUND\_HALF\_EVEN) and
\textbf{must} handle the following edge cases:

\begin{enumerate}
  \item Negative subtotals (credit notes): return negative VAT amount
  \item Zero rate: return VAT amount of exactly 0.00
  \item Rate $>$ 100\%: reject with error code \texttt{INVALID\_RATE}
\end{enumerate}

\subsection{Activity lookup}

\begin{lstlisting}[language=Python,basicstyle=\ttfamily\small,caption={Activity lookup API contract.}]
def lookup_activity_code_treatment(
    activity_code: str,
    country: str  # ISO 3166-1 alpha-2
) -> TreatmentResult:
    """
    Determine VAT treatment from business activity code.
    
    Returns:
        TreatmentResult with fields:
        - vat_category: str  # ZR, EX, FI, etc.
        - rate_percent: Optional[Decimal]
        - source: str  # regulation reference
    """
\end{lstlisting}

\begin{table}[htbp]
\centering
\caption{Activity code to VAT treatment mapping (Egypt, partial).}
\label{tab:activity-lookup}
\footnotesize
\begin{tabularx}{\linewidth}{@{}l l l X@{}}
\toprule
\textbf{Code} & \textbf{Description} & \textbf{Treatment} & \textbf{Source} \\ \midrule
1811 & Printing services & FI (14\%) & ETA Circular 2019/1 \\
8890 & Professional services & EX (exempt) & ETA Circular 2020/3 \\
4911 & Passenger rail transport & ZR (0\%) & VAT Law Art. 23 \\
6411 & Banking services & EX (exempt) & VAT Law Art. 29 \\
\bottomrule
\end{tabularx}
\end{table}

\section{JSON-Logic expression specification}

The rules engine evaluates predicates expressed in JSON-Logic format
(\url{https://jsonlogic.com}). This section specifies the supported
operators and variable paths.

\subsection{Supported operators}

\begin{table}[htbp]
\centering
\caption{JSON-Logic operators supported by the rules engine.}
\label{tab:jsonlogic-ops}
\footnotesize
\begin{tabularx}{\linewidth}{@{}l l X@{}}
\toprule
\textbf{Category} & \textbf{Operators} & \textbf{Notes} \\ \midrule
Comparison & \texttt{==}, \texttt{===}, \texttt{!=}, \texttt{!==} & Strict equality recommended \\
           & \texttt{<}, \texttt{<=}, \texttt{>}, \texttt{>=} & Numeric comparison \\
Logic      & \texttt{and}, \texttt{or}, \texttt{!}, \texttt{!!} & Boolean operations \\
           & \texttt{if} & Ternary conditional \\
String     & \texttt{in}, \texttt{cat}, \texttt{substr} & String operations \\
Array      & \texttt{in}, \texttt{merge}, \texttt{map}, \texttt{filter} & Array operations \\
Arithmetic & \texttt{+}, \texttt{-}, \texttt{*}, \texttt{/}, \texttt{\%} & Basic arithmetic \\
Data       & \texttt{var}, \texttt{missing}, \texttt{missing\_some} & Variable access \\
\bottomrule
\end{tabularx}
\end{table}

\subsection{Variable paths}

Variable paths use dot notation to access nested fields in the invoice record.
The root object is the validated invoice conforming to
\texttt{invoice.schema.json}.

\begin{lstlisting}[basicstyle=\ttfamily\small,caption={Example variable paths for invoice data.}]
totals.gross           # Invoice gross amount
totals.vat_rate_percent  # VAT rate as percentage
seller.tax_id          # Seller VAT registration number
seller.country         # Seller country code
buyer.country          # Buyer country code
metadata.issue_date    # Invoice issue date
line_items.0.description  # First line item description
\end{lstlisting}

\subsection{Named helper invocation}

When a predicate requires complex logic beyond JSON-Logic capabilities,
the engine invokes a named helper. The syntax is:

\begin{lstlisting}[basicstyle=\ttfamily\small,caption={Named helper invocation in rule predicate.}]
{
  "predicate": {
    "type": "named_helper",
    "name": "check_vat_registration_validity",
    "args": {
      "tax_id": {"var": "seller.tax_id"},
      "country": {"var": "seller.country"}
    }
  }
}
\end{lstlisting}

\section{Error handling specification}

The module \textbf{must} handle errors gracefully without crashing the
parent invoice processing pipeline.

\subsection{Error codes}

\begin{table}[htbp]
\centering
\caption{VAT module error codes.}
\label{tab:vat-errors}
\footnotesize
\begin{tabularx}{\linewidth}{@{}l l X@{}}
\toprule
\textbf{Code} & \textbf{Severity} & \textbf{Description} \\ \midrule
\texttt{VAT\_E001} & ERROR & Invalid VAT registration format \\
\texttt{VAT\_E002} & ERROR & Unknown country code \\
\texttt{VAT\_E003} & ERROR & Invalid VAT rate (negative or $>$100\%) \\
\texttt{VAT\_E004} & ERROR & JSON-Logic parse failure \\
\texttt{VAT\_W001} & WARNING & Authority API unavailable (fallback to format check) \\
\texttt{VAT\_W002} & WARNING & Activity code not in lookup table \\
\texttt{VAT\_I001} & INFO & Rate changed since last validation \\
\bottomrule
\end{tabularx}
\end{table}

\subsection{Fallback behaviour}

When an error occurs:

\begin{enumerate}
  \item Log the error with full context (invoice ID, field values, stack trace)
  \item Set the rule result to \textsc{hold} (never silently pass)
  \item Populate \texttt{checklist[].error\_code} with the appropriate code
  \item Continue processing remaining rules (fail-open for diagnostics)
  \item Return the partial checklist to the caller
\end{enumerate}

\section{Performance requirements}

\begin{table}[htbp]
\centering
\caption{VAT module performance requirements.}
\label{tab:vat-perf}
\footnotesize
\begin{tabularx}{\linewidth}{@{}l r X@{}}
\toprule
\textbf{Metric} & \textbf{Target} & \textbf{Measurement} \\ \midrule
p95 latency per invoice & $<$ 100 ms & Excluding authority API calls \\
p95 latency with API    & $<$ 2000 ms & Including authority validation \\
Memory per invoice      & $<$ 10 MB & Peak allocation during processing \\
Throughput              & $>$ 100/sec & Single-threaded, no API calls \\
\bottomrule
\end{tabularx}
\end{table}

\section{Implementation checklist}

Developers \textbf{must} verify the following before marking the module
complete:

\begin{itemize}
  \item[\texttt{[ ]}] All four helper functions implemented with correct signatures
  \item[\texttt{[ ]}] Unit tests achieve $\geq$90\% line coverage
  \item[\texttt{[ ]}] All country-specific regex patterns validated against test data
  \item[\texttt{[ ]}] JSON-Logic parser handles all operators in Table~\ref{tab:jsonlogic-ops}
  \item[\texttt{[ ]}] Error codes match Table~\ref{tab:vat-errors}
  \item[\texttt{[ ]}] Performance meets targets in Table~\ref{tab:vat-perf}
  \item[\texttt{[ ]}] Integration test with real invoice records from each country
  \item[\texttt{[ ]}] Security review completed (input validation, injection prevention)
\end{itemize}
